\ssr{ВВЕДЕНИЕ}

Современная стоматология активно использует цифровые методы диагностики. Одним из основных источников диагностической информации являются рентгеновские, панорамные и томографические снимки челюсти пациента. Такие снимки используются врачами для выявления патологий, планирования лечения, контроля динамики состояния пациента и формирования медицинских заключений.

Рост объема цифровых медицинских данных приводит к необходимости централизованного хранения снимков и связанных с ними сведений: информации о пациенте, типе исследования, оборудовании, клинике, враче и поставленном диагнозе. Использование разрозненных файловых хранилищ и ручного учета усложняет поиск данных, повышает риск потери информации и затрудняет контроль доступа к медицинским сведениям.

Актуальность работы заключается в необходимости разработки специализированной базы данных и приложения доступа, ориентированных на хранение и обработку стоматологических снимков челюсти человека. Такая система должна обеспечивать целостность данных, поддержку связей между медицинскими сущностями, разграничение доступа пользователей и возможность повышения скорости получения часто используемой информации за счет внешнего кэширования.

Целью курсовой работы является проектирование и разработка базы данных для хранения и обработки стоматологических снимков челюсти человека, а также приложения доступа к этой базе данных.

Для достижения поставленной цели необходимо решить следующие задачи:
\begin{itemize}
	\item проанализировать предметную область хранения и обработки стоматологических снимков;
	\item выполнить сравнительный анализ существующих решений;
	\item сформулировать требования к разрабатываемой базе данных и приложению доступа;
	\item формализовать данные, подлежащие хранению в базе данных;
	\item спроектировать реляционную схему базы данных и проверить ее соответствие третьей нормальной форме;
	\item спроектировать ролевую модель на уровне базы данных;
	\item реализовать базу данных и приложение доступа к данным;
	\item реализовать внешний механизм кэширования с использованием Redis;
	\item провести тестирование разработанного функционала;
	\item исследовать влияние применения индексов на время выполнения запросов к базе данных.
\end{itemize}

Объектом разработки является информационная система стоматологической клиники, обеспечивающая хранение и обработку цифровых снимков челюсти человека. Предметом разработки является реляционная база данных и приложение доступа к ней.

\clearpage
